% Generated draft: requires mathematical and source review.
\documentclass[11pt]{article}
\usepackage[margin=25mm]{geometry}
\usepackage{fontspec}
\setmainfont{Latin Modern Roman}
\usepackage{amsmath,amssymb,mathtools,unicode-math}
\setmathfont{Latin Modern Math}
\usepackage{enumitem,needspace,xurl,hyperref}
\setlength{\parindent}{0pt}
\setlength{\parskip}{5pt}
\setlength{\emergencystretch}{3em}
\newcommand{\N}{\mathbb{N}}
\newcommand{\Z}{\mathbb{Z}}
\newcommand{\Q}{\mathbb{Q}}
\newcommand{\R}{\mathbb{R}}
\newcommand{\C}{\mathbb{C}}
\newcommand{\F}{\mathbb{F}}
\newcommand{\A}{\mathbb{A}}
\newcommand{\PP}{\mathbb{P}}
\newcommand{\E}{\mathbb{E}}
\newcommand{\eps}{\varepsilon}
\newcommand{\dd}{\,\mathrm{d}}
\begin{document}
\section*{UnsolvedMath ID: 20000995}
\textbf{Source problem code:} catalogue entry
\subsection{Definitions and mathematical statement}
Persi Diaconis's source question asks what a limiting-shape theory for generalized internal diffusion-limited aggregation (IDLA) should look like, referring to [22], [4], and [21]. The record does not define the generalization, graph, walk kernel, source distribution, scaling, topology of convergence, or those references. It is a broad theory-building question.

The supplied aid treats a precise one-dimensional special case. Fix \(p+q=1\) with \(p>q>0\). Starting with \(A_0=\varnothing\subset\Z\), launch particles successively from zero. Each particle uses an independent nearest-neighbor walk with right-step probability \(p\), left-step probability \(q\), and settles at the first site outside the existing aggregate. Thus \(|A_N|=N\), and for \(N\ge1\),
\[
 A_N=[-L_N,R_N]\cap\Z,\qquad L_N,R_N\ge0,
 \qquad L_N+R_N+1=N.
\]
The upstream occupation length is \(L_N=-\min A_N\), opposite the positive drift. The aid reports the almost-sure boundary law
\[
 \frac{L_N}{\log_{p/q}N}\longrightarrow1\qquad(N\to\infty).
\]
Its proof is described as sandwiching changing-boundary waiting times by biased gambler's-ruin probabilities. Two explicit finite-\(N\) upper/lower tail inequalities are said to appear in the aid's Theorem 5.1, but their statements and constants are not supplied, so they cannot be reproduced exactly here. The law assumes fixed strictly positive bias and does not cover the symmetric case or a bias varying with \(N\).

The aid separately gives a conditional scaling principle for higher-dimensional kernels. Centered kernels, under suitable nondegeneracy and shape-theorem hypotheses, are associated with elliptic scaling of order \(N^{1/d}\) in all directions. With fixed nonzero drift and \(d\ge2\), it reports parabolic scales
\[
 N^{1/(d+1)}\quad\text{transverse to the drift},\qquad
 N^{2/(d+1)}\quad\text{along the drift}.
\]
The generator balance is between longitudinal first-order drift and transverse second-order diffusion, and the product of one longitudinal and \(d-1\) transverse scales has order \(N\). This expansion motivates candidate exponents; it is not a shape theorem for arbitrary generalized IDLA. Moment assumptions, ellipticity/nondegeneracy, geometry, limiting free-boundary problem, and convergence estimates are not supplied. In dimension one the reported logarithmic upstream boundary is a finer feature than the order-\(N\) downstream scale. The broad source question remains only partially addressed.
\subsection{Short English statement}
The source asks for a general IDLA shape theory. The supplied aid gives a sharp logarithmic upstream law for a fixed biased walk on \(\Z\) and distinguishes centered elliptic scaling from conditional drift-driven parabolic scaling, without supplying a universal shape theorem.
\subsection{Sources}
\begin{itemize}
\item Local UnsolvedMath registry, permanent ID \(20000995\), source code \texttt{catalogue entry}; \texttt{aim-combinatorics-notes.json}, \texttt{notes[119]}, tag \texttt{problem}.
\item AIM workshop ``Generalizations of chip-firing and the critical group,'' Problem 30, Persi Diaconis: \url{https://aimath.org/pastworkshops/chipfiringproblems.pdf}. Statement status \texttt{exact}. Original reference labels [22], [4], and [21] are supplied without their bibliographic entries. No PDF was accessed for this draft.
\item The supplied aid refers to its Theorem 5.1 for two finite-size tail inequalities; no theorem text, artifact, or authoritative external citation is included. No original remarks or literature field are supplied.
\item Research attempt 1: \texttt{partial\_result}, status \texttt{partially\_solved}, classification \texttt{PARTIAL-PROGRESS}, rigor label \texttt{proved}, novelty confidence low. Created and updated 14 August 2026. This machine-generated aid is not peer reviewed; independently verify before citation.
\item Ulam AI / UnsolvedMath Contributors; CC BY 4.0.
\end{itemize}
\end{document}
